\documentclass[12pt]{article}
\usepackage[T1]{fontenc}
\usepackage[utf8]{inputenc}
\usepackage{lmodern}
\usepackage{textcomp}
\usepackage{enumitem}
\usepackage{geometry}
\geometry{margin=1in}
\begin{document}
\begin{center}
Answer Key - Astronomy - Graduate Level Assessment 19 \
\small (Answers are ordered exactly as in the question paper.)
\end{center}

\section*{Multiple Choice}
\begin{enumerate}[label=\arabic*.]
\item \textbf{Answer:} A
\\ \textit{Options:} A) Its rotation is too slow.; B) It has too thick an atmosphere.; C) It is too large.; D) It does not have a metallic core.
\\ \textit{Note:} Venus has a slow rotation period of about 243 Earth days, which is insufficient to generate the dynamo effect necessary for creating a significant magnetic field, unlike Earth, which rotates more rapidly.
\item \textbf{Answer:} A
\\ \textit{Options:} A) We don't know for sure.; B) aliens (panspermia).; C) comets.; D) the Earth's oceans.
\\ \textit{Note:} The correct answer is supported by the fact that despite various hypotheses regarding the origin of life on Earth, including abiogenesis and panspermia, there is currently no definitive scientific consensus or empirical evidence to conclusively determine how life first emerged.
\item \textbf{Answer:} B
\\ \textit{Options:} A) nuclear fusion in the core; B) by contracting changing gravitational potential energy into thermal energy; C) internal friction due to its high rotation rate; D) chemical processes
\\ \textit{Note:} Astronomers believe Jupiter generates its internal heat through the process of gravitational contraction, where the planet's immense mass increases pressure as it contracts, converting gravitational potential energy into thermal energy, which contributes to its overall heat.
\item \textbf{Answer:} D
\\ \textit{Options:} A) Schwarzschild radiation; B) Planck radiation; C) Kolmogorov radiation; D) Hawking radiation
\\ \textit{Note:} Hawking radiation is the theoretical prediction that black holes emit radiation due to quantum effects near the event horizon, leading to their gradual evaporation over time.
\item \textbf{Answer:} C
\\ \textit{Options:} A) Mars; B) Venus; C) Earth; D) Mars and Earth
\\ \textit{Note:} Earth is the only planet in the solar system known to have active plate tectonics, characterized by the movement of tectonic plates that shape its surface and contribute to geological phenomena such as earthquakes and volcanic activity.
\end{enumerate}

\section*{True / False}
\begin{enumerate}[label=\arabic*.]
\setcounter{enumi}{5}
\item \textbf{Answer:} True
\\ \textit{Options:} A) True; B) False
\\ \textit{Note:} The statement is true because Phobos and Deimos are indeed the names of Mars' two natural satellites, which were discovered in 1877 by American astronomer Asaph Hall.
\item \textbf{Answer:} True
\\ \textit{Options:} A) True; B) False
\\ \textit{Note:} The thin atmosphere of Mars, composed mainly of carbon dioxide but with a surface pressure less than 1\% of Earth's, is insufficient to effectively retain heat, thereby significantly diminishing the greenhouse effect's capacity to warm the planet.
\item \textbf{Answer:} True
\\ \textit{Options:} A) True; B) False
\\ \textit{Note:} Many organisms, such as anaerobic bacteria and certain extremophiles, can survive and thrive in environments devoid of oxygen by utilizing alternative metabolic pathways, demonstrating that oxygen is not universally essential for life.
\item \textbf{Answer:} False
\\ \textit{Options:} A) True; B) False
\\ \textit{Note:} The Andromeda Galaxy is actually located about 2.537 million light years away from Earth, which makes the statement inaccurate.
\item \textbf{Answer:} False
\\ \textit{Options:} A) True; B) False
\\ \textit{Note:} While Jupiter does have equatorial wind speeds that exceed 900 miles per hour, Saturn's equatorial winds are significantly lower, averaging around 1,100 kilometers per hour (approximately 680 miles per hour), making the statement false.
\end{enumerate}

\section*{Long Form Response}
\begin{enumerate}[label=\arabic*.]
\setcounter{enumi}{10}
\item \textbf{Answer:} The Mars Exploration Rovers Opportunity and Spirit face significant operational limitations that hinder their research capabilities. Opportunity's ability to maneuver its arm may be compromised, limiting its capacity to conduct detailed geological analyses and collect samples, which are crucial for understanding Mars' past environment. Meanwhile, Spirit is at risk of insufficient solar power during the approaching winter months, potentially reducing its operational time and effectiveness in gathering data. These challenges not only restrict the rovers' immediate scientific objectives but also raise concerns about the long-term sustainability of Mars exploration efforts. The implications for ongoing research are profound, as they may result in missed opportunities to uncover critical insights into Martian geology and climate, which could inform future missions.
\\ \textit{Note:} correct because it identifies specific operational limitations of the rovers-Opportunity's compromised arm movement affecting geological analysis and Spirit's risk of reduced solar power-highlighting their implications for ongoing Mars research.
\item \textbf{Answer:} The Arecibo Telescope, operational from 1963 until its collapse in 2020, played a pivotal role in the advancement of radio astronomy due to its unique features, including its large 305-meter dish, which enabled high sensitivity and resolution in detecting radio waves from celestial objects. Its contributions to our understanding of the universe are substantial, as it facilitated groundbreaking research in areas such as pulsar timing, planetary radar observations, and the study of cosmic phenomena like gravitational waves. The telescope's ability to transmit and receive radio signals allowed for detailed studies of the solar system, including mapping the surfaces of planets and moons, and it was instrumental in the discovery of the first binary pulsar, which provided crucial evidence for the existence of gravitational radiation. Arecibo's legacy also includes its role in the Search for Extraterrestrial Intelligence (SETI), as it was used to broadcast messages into space, symbolizing humanity's quest to connect with potential extraterrestrial civilizations. Ultimately, the Arecibo Telescope not only advanced scientific knowledge but also inspired public interest in astronomy and the search for life beyond Earth.
\\ \textit{Note:} correct because it accurately identifies the Arecibo Telescope's significant features, such as its large dish size, and details its contributions to key areas of radio astronomy, thereby underscoring its pivotal role in enhancing our understanding of the universe.
\end{enumerate}

\vfill
\noindent\textit{End of Answer Key}
\end{document}